\input{preamble}
\hypersetup{pdftitle={Reductions toward the Riemann hypothesis},
  pdfauthor={GRH distillation project}}
\fancyhead[L]{\small Reductions toward the Riemann hypothesis}
\fancyhead[R]{\small RH manuscript}
\title{Reductions toward the Riemann hypothesis\\
  \large A short analytic endpoint and a compatible arithmetic target}
\author{GRH distillation project}
\date{November 2026}
\begin{document}
\maketitle
\begin{abstract}
The Riemann hypothesis follows from an eventual upper bound of order
$x^{1/2+\epsilon}$ for a corrected smoothed prime sum. We give the full
Mellin-transform and positivity proof, and state a stronger sufficient
condition using weighted positive-part averages on arbitrarily large
intervals. A direct arithmetic comparison uses positivity of the full
primitive-character kernel, including ramified primes. Averaging after
this comparison leaves four quadratic characters and requires no symmetry
of the individual weights. The remaining geometric estimates are organized
around one component mean volume, short integral sections, and exact
aggregate energies of the same arithmetic maps. These yield conditional
degree and slope bounds with explicit losses. The required compatible
arithmetic construction remains open; this manuscript is a collection of
proved reductions and precisely stated sufficient conditions, not a proof
of the Riemann hypothesis.
\end{abstract}
\setcounter{tocdepth}{1}
\tableofcontents
\clearpage
\input{rh/sections/01-introduction}
\input{rh/sections/02-analytic-endpoint}
\input{rh/sections/03-direct-transfer}
\input{rh/sections/04-geometric-degree}
\input{rh/sections/05-theta-energy}
\input{rh/sections/06-compatible-budget}
\input{rh/sections/07-references}
\end{document}
